% © 2026 Ing. David Jaroš — CC BY-NC-ND 4.0
%
% This work is licensed under a Creative Commons Attribution-NonCommercial-NoDerivatives
% 4.0 International License (CC BY-NC-ND 4.0).
%
% License History: Earlier drafts (up to v0.3) were released under CC BY 4.0.
% From v0.4 onward, all material is released under CC BY-NC-ND 4.0 to protect
% the integrity of the theoretical work during ongoing academic development.
%
% See LICENSE.md for full license text.
%
% File: canonical/alpha/veff_corrected_statement.tex
% Purpose: Corrected formal statement of V_eff(n) and its stationary-point condition,
%          following the V_eff sanity audit (reports/alpha_veff_sanity_audit.md).
%          Replaces the erroneous closed-form n*(B) = e^{(2-B)/B}·e previously
%          stated in alpha_equation_matrix.tex.
% Status: 2026-05-02

% UBT-AI-PROVENANCE-BEGIN
% schema: ubt-ai-provenance/v1
% tier: B_machine_verified
% ai_assistance: disclosed
% human_review: machine-verification
% editorial_responsibility: Ing. David Jaroš
% policy: ../../AI_PROVENANCE.md
% notice: Machine-verified against named sources or verifiers; individual attestation is not claimed.
% UBT-AI-PROVENANCE-END

\ifdefined\INCLUDEMODE
\else
  \documentclass[12pt,a4paper]{article}
  \usepackage[a4paper, margin=2.5cm]{geometry}
  \usepackage{amsmath, amssymb, amsthm}
  \usepackage{mathtools}
  \usepackage{hyperref}
  \usepackage{xcolor}
  \usepackage{booktabs}
  \usepackage{longtable}
  \usepackage{mdframed}
  \usepackage{tcolorbox}
  \tcbuselibrary{skins,breakable}
  \usepackage{array}

  \newtheorem{theorem}{Theorem}[section]
  \newtheorem{lemma}[theorem]{Lemma}
  \newtheorem{proposition}[theorem]{Proposition}
  \newtheorem{corollary}[theorem]{Corollary}
  \theoremstyle{definition}
  \newtheorem{definition}[theorem]{Definition}
  \theoremstyle{remark}
  \newtheorem{remark}[theorem]{Remark}

  \newmdenv[backgroundcolor=yellow!20, linecolor=orange, linewidth=1.5pt]{gapbox}
  \newmdenv[backgroundcolor=green!15, linecolor=green!60!black, linewidth=1.5pt]{resultbox}
  \newmdenv[backgroundcolor=red!8,   linecolor=red!60,         linewidth=1.5pt]{warnbox}

  \newcommand{\BB}{\mathbb{B}}
  \newcommand{\CC}{\mathbb{C}}
  \newcommand{\RR}{\mathbb{R}}
  \newcommand{\HH}{\mathbb{H}}
  \newcommand{\ZZ}{\mathbb{Z}}
  \newcommand{\Tr}{\mathrm{Tr}}

  \newcommand{\PROVED}{\textbf{[Proved]}}
  \newcommand{\OPEN}{\textbf{[Open]}}
  \newcommand{\COND}{\textbf{[COND]}}
  \newcommand{\MC}{\textbf{[MC]}}
  \newcommand{\Lone}{\textbf{[L1]}}
  \newcommand{\Lzero}{\textbf{[L0]}}
  \newcommand{\PHENOM}{\textbf{[PHENOM]}}

  \title{\textbf{Corrected Statement of $V_{\mathrm{eff}}(n)$ and the\\
         $n^* = 137$ Stationarity Condition}\\
         \large Following the V\_eff Sanity Audit of 2026-05-02}
  \author{Ing.\ David Jaro\v{s}}
  \date{May 2026}

  \begin{document}
  \maketitle
\fi

%──────────────────────────────────────────────────────────────────────────────
\section{Purpose and Scope}
\label{sec:vc:purpose}
%──────────────────────────────────────────────────────────────────────────────

This document supersedes the stationarity formula for $V_{\mathrm{eff}}(n)$
previously stated in \texttt{canonical/alpha/alpha\_equation\_matrix.tex},
which contained an error in the closed-form expression for $n^*(B)$.
It provides:

\begin{enumerate}
  \item The correct definition of $V_{\mathrm{eff}}(n)$ and all variants found in the repo.
  \item The correct stationarity condition (transcendental; no closed form).
  \item The exact value of~$B$ required for $n^* = 137$.
  \item A classification of every~$B$ value in the repository.
  \item The decision on whether the $n^* = 137$ claim is mathematically valid.
\end{enumerate}

\begin{warnbox}
\textbf{Hard rule (unchanged)}: No constant in this document is chosen to
match~$\alpha$ or~$137$.  All constants are independently derived or explicitly
labelled with their proof status.
\end{warnbox}

%──────────────────────────────────────────────────────────────────────────────
\section{Catalogue of $V_{\mathrm{eff}}$ Variants}
\label{sec:vc:variants}
%──────────────────────────────────────────────────────────────────────────────

A systematic search of the repository identified the following distinct
effective-potential formulas.

\subsection{V1 — Primary Canonical Form}
\label{sec:vc:V1}

\begin{equation}
  \boxed{V_{\mathrm{eff}}(n) \;=\; n^2 \;-\; B\,n\ln n,}
  \qquad n \in \ZZ_{>0}.
  \label{eq:vc:Veff_canonical}
\end{equation}

\noindent\textbf{Sources}:
\texttt{canonical/alpha/alpha\_equation\_matrix.tex},
\texttt{canonical/alpha/alpha\_best\_route.tex~\S6},
\texttt{canonical/alpha/neff\_32\_alpha\_route.tex~\S G5},
\texttt{canonical/appendices/appendix\_alpha\_geometry.tex~\S3},
\texttt{docs/papers/papers/generated/ubt\_action\_and\_alpha.tex}.

\noindent\textbf{Origin of $A=1$}: The classical Kaluza-Klein winding energy
$\hbar^2n^2/(2mR_\psi^2)$ equals $n^2$ in natural units at $R_\psi = 1$.
This is a choice of units, not a free parameter.

\textbf{Status}: \PROVED~\Lone\ (functional form); coefficient $B$ is an open gap.

\subsection{V2 — Legacy Code Variant ($A$ explicit)}

\begin{equation}
  V_{\mathrm{eff}}(n;\,A,B) \;=\; A\,n^2 \;-\; B\,n\ln n,
  \quad A = 1\;\text{(default)},
  \label{eq:vc:Veff_legacy}
\end{equation}
with $B = 20.3$ in the earliest script
(\texttt{ARCHIVE/archive\_legacy/tex/emergent\_alpha\_calculations.tex}).

\begin{warnbox}
The script with $B = 20.3$ restricts the prime search to $p \in [100, 200]$,
creating the false appearance that $B = 20.3$ selects $n = 137$.
The \emph{global} prime minimiser of~\eqref{eq:vc:Veff_legacy} for $B = 20.3$
is $p = 47$, not $p = 137$.
This result must not be cited as evidence for the prime attractor.
\end{warnbox}

\subsection{V3 — Non-Hermitian Mass Term (Unrelated to $\alpha$)}

\begin{equation}
  V_{\mathrm{eff}} = (m^2 + i\gamma)\,\mathrm{Tr}[\Theta^\dagger\Theta],
  \quad \gamma \in \RR.
  \label{eq:vc:Veff_NH}
\end{equation}
Source: \texttt{canonical/symmetry/step3\_breaking\_catalogue.tex}.
This potential describes open-system $\mathcal{PT}$-symmetric dynamics.
It plays \textbf{no role in the $\alpha$ derivation}.

\subsection{V4 — Erroneous Form (Historical Bug)}

\begin{equation}
  V(n) \;=\; n^2 \;-\; B\ln n \qquad \text{(missing factor }n\text{)}.
  \label{eq:vc:Veff_wrong}
\end{equation}
This variant appears nowhere as an explicit formula, but it is the only form
consistent with the stationarity claim $n^* = \sqrt{B/2}$ quoted in
\texttt{ALPHA\_PROGRESS\_REPORT.md~\S2.4}.  For $B \approx 46.3$,
$\sqrt{B/2} \approx 4.8$, which is clearly not 137.
This variant is a \textbf{transcription error} and should be disregarded.

%──────────────────────────────────────────────────────────────────────────────
\section{Corrected Stationarity Condition}
\label{sec:vc:stationary}
%──────────────────────────────────────────────────────────────────────────────

\begin{theorem}[Stationarity of $V_{\mathrm{eff}}$, corrected]
\label{thm:vc:stationary}
The continuous minimum of $V_{\mathrm{eff}}(n) = n^2 - B\,n\ln n$ satisfies
the \textbf{transcendental equation}:
\begin{equation}
  \frac{\partial V_{\mathrm{eff}}}{\partial n}\bigg|_{n^*}
  = 2n^* - B\bigl(\ln n^* + 1\bigr) = 0,
  \label{eq:vc:stationary}
\end{equation}
equivalently
\begin{equation}
  \boxed{2n^* = B\,\bigl(\ln n^* + 1\bigr).}
  \label{eq:vc:stationary_box}
\end{equation}
This equation has \textbf{no elementary closed form} for $n^*(B)$.
\end{theorem}

\begin{warnbox}
\textbf{Correction of previous error}:
The formula $n^*(B) = e^{(2-B)/B}\cdot e$ stated in
\texttt{alpha\_equation\_matrix.tex} (Theorem~\ref{thm:vc:stationary},
§\,Shared Foundation) is \textbf{incorrect}.

Algebraically: $e^{(2-B)/B}\cdot e = e^{(2-B)/B+1} = e^{2/B}$.
For $B = 46.298$: $e^{2/B} \approx 1.044$, not $137$.

The correct relation is the transcendental Eq.~\eqref{eq:vc:stationary_box}.
\end{warnbox}

\begin{corollary}[No closed form for $n^*(B)$]
The function $n^*(B)$ defined by~\eqref{eq:vc:stationary_box} is the solution
of a transcendental equation and is evaluated numerically for each~$B$.
Selected values:
\begin{center}
\begin{tabular}{lll}
\toprule
$B$ value & $n^*_{\mathrm{continuous}}$ & Note \\
\midrule
$8\pi \approx 25.133$ & $65.0$ & $B_0$, proved [L1] \\
$N_{\mathrm{eff}}^{3/2} \approx 41.569$ & $120.4$ & $B_{\mathrm{base}}$, conjectural [MC] \\
$\mu(\Gamma_0(137))/3 = 46.000$ & $136.0$ & exact math, 0.64\% below $B_{\mathrm{phenom}}$ \\
$46.284$ & $137.000$ & exact constraint for $n^* = 137$ \\
$46.298$ & $137.05$ & $B_{\mathrm{phenom}}$ (phenomenological) \\
\bottomrule
\end{tabular}
\end{center}
\end{corollary}

\begin{corollary}[Second derivative confirms minimum]
\label{cor:vc:2nd}
\[
  \frac{\partial^2 V_{\mathrm{eff}}}{\partial n^2}\bigg|_{n^*}
  = 2 - \frac{B}{n^*}
  \;>\; 0
  \quad\Longleftrightarrow\quad
  n^* > \frac{B}{2}.
\]
For $B \approx 46.3$ and $n^* = 137$: $137 > 23.15$\,\checkmark.\
The stationary point is indeed a \textbf{minimum}.
\end{corollary}

%──────────────────────────────────────────────────────────────────────────────
\section{Required Coefficient for $n^* = 137$}
\label{sec:vc:B_required}
%──────────────────────────────────────────────────────────────────────────────

\begin{proposition}[Exact $B$ for $n^* = 137$]
\label{prop:vc:B_req}
Substituting $n^* = 137$ into the stationarity condition~\eqref{eq:vc:stationary_box}:
\begin{equation}
  B_{\mathrm{required}}
  = \frac{2 \times 137}{\ln 137 + 1}
  = \frac{274}{4.919981 + 1}
  = \frac{274}{5.919981}
  \approx 46.284.
  \label{eq:vc:B_req}
\end{equation}
The phenomenological value $B_{\mathrm{phenom}} = 46.298$ in the repository
differs from $B_{\mathrm{required}}$ by $0.030\%$, consistent with rounding.
\end{proposition}

\begin{corollary}[Prime minimizer]
For any $B$ in the range $[46.00,\,46.50]$, the prime minimiser of
$V_{\mathrm{eff}}(p)$ over all primes $p$ is $p = 137$:
\begin{equation}
  V_{\mathrm{eff}}(137) < V_{\mathrm{eff}}(p)
  \quad\text{for all primes } p \neq 137,\quad B \in [46.00,\,46.50].
  \label{eq:vc:prime_min}
\end{equation}
The next lowest prime is $p = 139$ with $V(139) - V(137) \approx +3.16$.
\end{corollary}

%──────────────────────────────────────────────────────────────────────────────
\section{Classification of All $B$ Values}
\label{sec:vc:B_class}
%──────────────────────────────────────────────────────────────────────────────

\begin{center}
\renewcommand{\arraystretch}{1.3}
\begin{tabular}{p{3.8cm} p{1.8cm} p{3.5cm} p{3.8cm}}
\toprule
\textbf{Symbol} & \textbf{Value} & \textbf{Classification} & \textbf{Proof status / source} \\
\midrule
$B_0 = 2\pi N_{\mathrm{eff}}/3$ & $8\pi\approx 25.13$ & DERIVED & \PROVED~\Lone; one-loop vacuum polarisation \\
$B_{\mathrm{base}} = N_{\mathrm{eff}}^{3/2}$ & $\approx 41.57$ & CONJECTURAL & \MC; conditional on $k_{\mathrm{KM}}=1$ (Gap G3-k) \\
$\mu(\Gamma_0(137))/3$ & $46.00$ & EXACT MATH & \PROVED~\Lzero\ (index formula); identification with $B$ is \OPEN \\
$B_{\mathrm{required}}$ & $46.284$ & EXACT CONSTRAINT & Unique value with $n^*_{\mathrm{continuous}}=137$; \PHENOM \\
$B_{\mathrm{phenom}}$ & $\approx 46.298$ & PHENOMENOLOGICAL & Back-solved from $n^*=137$; circular unless Gap G137-B is formally resolved \\
$B_{\mathrm{full}} = B_{\mathrm{base}}\cdot R$ & $\approx 46.31$ & OPEN & $R\approx 1.114$; no non-circular derivation \\
\bottomrule
\end{tabular}
\end{center}

\begin{gapbox}
\textbf{Gap G137-B (primary blocker)}:
Derive $B \approx 46.284$ from $S[\Theta]$ without using $\alpha$ as input.
Equivalently: derive the correction factor $R = B_{\mathrm{full}}/B_{\mathrm{base}}
\approx 1.1134$ from first principles.

Until this gap is closed, the claim $n^*(B) = 137$ is \textbf{CONDITIONAL}
on $B = B_{\mathrm{phenom}}$, not a zero-parameter result.
\end{gapbox}

%──────────────────────────────────────────────────────────────────────────────
\section{Audit Verdict and Decision}
\label{sec:vc:decision}
%──────────────────────────────────────────────────────────────────────────────

\begin{tcolorbox}[colback=blue!4!white,colframe=blue!50!black,title=Audit Verdict]
\begin{enumerate}
  \item \textbf{V\_eff functional form is correct.}
        $V_{\mathrm{eff}}(n) = n^2 - B\,n\ln n$ is the unique form consistent
        with the one-loop Coleman-Weinberg / heat-kernel calculation on $S^1_\psi$.

  \item \textbf{The $n^* = 137$ extremum claim is mathematically valid}
        for $B \approx 46.284$--$46.298$.  The stationarity condition is correct
        and the prime minimiser is 137 in this $B$-range. \Lone~(given $B$).

  \item \textbf{The closed-form $n^*(B) = e^{(2-B)/B}\cdot e$ in
        \texttt{alpha\_equation\_matrix.tex} is an error.}  The correct relation
        is the transcendental equation~\eqref{eq:vc:stationary_box}.
        (See \S\ref{sec:vc:stationary}.)

  \item \textbf{$B \approx 46.3$ is phenomenological.}
        It is obtained by back-substituting $n^* = 137$ into the stationarity
        condition.  It is circular unless Gap G137-B is independently closed.

  \item \textbf{$B_{\mathrm{phenom}}$ does not require a hidden scale factor beyond
        the well-understood missing factor $R \approx 1.113$}, but that factor
        itself is open.  The closest derived value is $\mu(\Gamma_0(137))/3 = 46.00$
        (0.64\% discrepancy), which is a structural corroboration, not a proof.
\end{enumerate}
\end{tcolorbox}

\subsection*{Next steps (from audit)}

\begin{enumerate}
  \item \textbf{Correct the formula in \texttt{alpha\_equation\_matrix.tex}}:
        replace $n^*(B) = e^{(2-B)/B}\cdot e$ with the transcendental
        condition~\eqref{eq:vc:stationary_box}.

  \item \textbf{Correct the formula in \texttt{ALPHA\_PROGRESS\_REPORT.md~\S2.4}}:
        replace $n^* = \sqrt{B/2}$ with the correct stationarity
        condition~\eqref{eq:vc:stationary_box}.

  \item \textbf{Continue attacking Gap G137-B}:
        primary attack via modular bootstrap on $\hat{Z}(\tau) = \vartheta_3^3(\tau)$.
        Secondary: explicit two-loop heat-kernel determinant on $T^2\times S^1_\psi$.

  \item \textbf{Do not promote $n^*(B_{\mathrm{phenom}}) = 137$ to PROVED}
        until Gap G137-B is formally resolved.
\end{enumerate}

\ifdefined\INCLUDEMODE
\else
  \end{document}
\fi
